\documentclass[11pt,a4paper]{article}

\usepackage[utf8]{inputenc}
\usepackage[T1]{fontenc}
\usepackage{amsmath,amssymb,amsfonts,amsthm,bm}
\usepackage{graphicx}
\usepackage{booktabs}
\usepackage{hyperref}
\hypersetup{hidelinks}
\usepackage{url}
\usepackage[margin=1in]{geometry}
\usepackage{enumitem}

\newtheorem{theorem}{Theorem}
\newtheorem{corollary}[theorem]{Corollary}

\title{Constraint Locality in Learned Syndrome Decoding\\for 3D Occupancy Grids}

\author{Brian Guarraci\\
  Galatheus\\
  \texttt{brian@galatheus.com}}

\date{May 2026}

\begin{document}
\maketitle

\begin{abstract}
We study learned reconstruction for underdetermined binary syndrome systems
arising from $32^3$ occupancy grids. A syndrome
$\mathbf{H}\mathbf{x}=\mathbf{s}\pmod 2$ specifies an exponentially large
coset, so reconstruction depends on a data-dependent learned proposal
conditioned on coarse occupancy, low-frequency spectral side information,
and syndrome-derived features. The central question is not whether a
syndrome can uniquely reconstruct a grid, but which constraint structures
help a learned proposal refine its prediction.

The main empirical finding is a constraint-locality effect. Hard algebraic
projection with fixed high-weight checks can degrade a plausible learned
prediction, while belief propagation on local overlapping surface checks can
improve it. We frame the result as a design rule for learned syndrome-based
occupancy coding: constraints should be local, overlapping, and targeted to
uncertain regions of the proposal. We do not claim a complete production codec
or state-of-the-art compression system; the contribution is the experimental
isolation of when algebraic constraints help or hurt learned reconstruction.
\end{abstract}

% ====================================================================
\section{Introduction}

Binary syndrome decoding---recovering $\mathbf{x} \in \{0,1\}^N$ from
$\mathbf{s} = \mathbf{H}\mathbf{x} \bmod 2$ given a parity check
matrix $\mathbf{H} \in \{0,1\}^{M \times N}$---is a fundamental
problem in coding theory~\cite{gallager}. When $M < N$, the system is
underdetermined: the syndrome constrains $\mathbf{x}$ to a coset of
size $2^{N-\text{rank}(\mathbf{H})}$, but does not uniquely determine
it. The decoder must select from this exponentially large set using
additional information.

In classical channel coding, the additional information is the channel
observation (soft bit likelihoods). In compressed
sensing~\cite{candes2006}, it is a sparsity prior. In this work, we
study a setting where the additional information is a \textbf{learned
spatial proposal}: a neural network trained on data from the target domain
(3D voxel grids of real objects) that predicts likely solutions from
spectral side-information and syndrome-derived features.

This setting arises naturally in 3D volumetric compression, where the
encoder computes a syndrome of a binary voxel grid and transmits it
alongside spectral coefficients. In this paper we keep the scope narrow:
we evaluate 3D occupancy-grid reconstruction and use compression as the
motivating rate-distortion setting. We do not rely on cross-domain analogies
for the core result.

This is the compression member of a four-paper program on learned recovery over
constraint systems. Its role is deliberately local: establish when parity
checks help or hurt a learned spatial proposal. Companion papers use that
constraint-locality lesson as framing for surface-code QEC, prove the
axial-parity gauge theorem, and test a structure-aware neural decoder for the
X-cube fracton code; none of those claims is needed for the empirical result
reported here.

The work is closest in spirit to syndrome source coding with decoder side
information~\cite{slepianwolf,wynerziv} and neural decoders for linear
codes~\cite{nachmani2016,nachmani2018}. We do not derive a new
rate-distortion theorem and we do not claim a state-of-the-art neural
compression system in the lineage of end-to-end learned image
codecs~\cite{balle2017,minnen2018}. Instead, we isolate a design question
that appears once a learned decoder is already strong: which parity-check
structures improve the learned proposal, and which structures damage it?

Our contributions:
\begin{enumerate}[nosep]
  \item \textbf{Diagnosis}: we evaluate how a fixed high-weight parity
    system interacts with a learned spatial proposal on $32^3$ ShapeNet
    occupancy grids (Section~\ref{sec:diagnosis}).
  \item \textbf{A constraint locality principle}: we compare hard
    algebraic projection against BP refinement with local overlapping
    surface checks and identify the structural conditions under which
    syndrome information improves the proposal (Section~\ref{sec:locality}).
    We support the observation with a simple check-weight analysis
    (Theorems~\ref{thm:utility}--\ref{thm:expected}).
  \item \textbf{Content-adaptive constraints}: we show that
    $\mathbf{H}$ can be derived from shared side-information
    (low-frequency spectral coefficients) available to both encoder and decoder,
    requiring no additional constraint-location signaling bits
    (Section~\ref{sec:adaptive}).
  \item \textbf{Density-stratified evidence}: we show that local BP helps
    in every occupancy quartile and helps most on the densest validation
    objects, while hard projection hurts in every quartile.
\end{enumerate}

% ====================================================================
\section{Problem Setting}

Let $\mathbf{x} \in \{0,1\}^N$ be a binary vector with spatial
structure (a $d^3$ voxel grid, $N = d^3$). The encoder computes:
\begin{itemize}[nosep]
  \item A syndrome $\mathbf{s} = \mathbf{H}\mathbf{x} \bmod 2 \in
    \{0,1\}^M$ from a parity check matrix $\mathbf{H}$.
  \item Spectral side-information $\mathbf{c}$ (e.g., low-frequency
    transform coefficients of $\mathbf{x}$).
\end{itemize}
The decoder observes $(\mathbf{s}, \mathbf{c})$ and must recover
$\mathbf{x}$.

With the corrected fixed constraint set, $M = 2{,}977$ and
$N = 32{,}768$; the syndrome eliminates only $2{,}977$ bits of
entropy, leaving $2^{29{,}791}$ valid solutions.
The spectral coefficients provide a soft spatial proposal
$P_\theta(x_i = 1 | \mathbf{c},\mathbf{s})$ through the learned decoder, but
the decoder's fidelity is bandwidth-dependent: the post-projection L2 result
is IoU $0.6947\pm0.0072$, while the L4 learned proposal before projection is
IoU $0.9370\pm0.0028$.
The question is: \textbf{how should the decoder combine the algebraic
constraints with the spatial proposal?}

\subsection{Three Decoding Operations}

\textbf{(A) Hard algebraic projection.}
Given a learned proposal $\tilde{\mathbf{x}}$, compute its residual
$\mathbf{r}=\mathbf{s}-\mathbf{H}\tilde{\mathbf{x}}$ over GF(2) and apply a
right-inverse correction
$\hat{\mathbf{x}}=\tilde{\mathbf{x}}+\mathbf{L}\mathbf{r}$, where
$\mathbf{H}\mathbf{L}=\mathbf{I}_M$. This satisfies the syndrome exactly but
chooses a correction by algebraic convenience rather than spatial coherence.
The canonical solution $\mathbf{L}\mathbf{s}$ is the special case
$\tilde{\mathbf{x}}=0$.

\textbf{(B) Learned proposal decoding.}
Train a neural network
$f_\theta(\mathbf{x}_0(\mathbf{s}),\mathbf{c},\mathbf{s}) \to
\hat{\mathbf{x}}$ on domain data, where $\mathbf{x}_0(\mathbf{s})$ is a
canonical syndrome solution. The network learns what spatially plausible
solutions look like, but its output is not guaranteed to satisfy the syndrome
exactly.

\textbf{(C) Proposal + constraint refinement.}
Use the learned proposal as an initial estimate, then refine via belief
propagation on the Tanner graph~\cite{tanner,kfl2001} of $\mathbf{H}$ with the
proposal's confidence as variable beliefs and $\mathbf{s}$ as check values.

We show that (A) fails in this setting, (B) is strong,
and (C) is better than (B) \emph{only when $\mathbf{H}$ has the right
structure}.

% ====================================================================
\section{Diagnosis: Why Hard Algebraic Projection Fails}
\label{sec:diagnosis}

The clean rsfix evaluation of operation~(A) uses $32^3$ ShapeNet voxel
grids~\cite{3dr2n2} with the corrected axis + plane + global parity checks
($M = 2{,}977$). After learned proposal decoding (operation~B, mean
IoU~$0.9370\pm0.0028$ at spectral level~4), we project onto the constraint
surface via the right inverse.

\begin{table}[htbp]
\centering
\caption{Clean hard algebraic projection diagnosis (200 ShapeNet validation
  samples, spectral level~4, mean $\pm$ SEM over seeds 42/43/44). The algebraic
  correction is usually wrong even though the corrected constraint set has
  only 1.6\% violated checks before projection.}
\label{tab:projection}
\begin{tabular}{lr}
\toprule
Metric & Value \\
\midrule
Constraint violations before projection & $47.6 \pm 4.5$ / 2{,}977 (1.6\%) \\
Bits flipped by projection & $52.5 \pm 5.2$ / 32{,}768 (0.16\%) \\
IoU before projection & $0.9370 \pm 0.0028$ \\
IoU after projection & $0.8782 \pm 0.0085$ (\textbf{$-0.0588 \pm 0.0058$}) \\
Mean $|\text{logit}|$ of flipped bits & $8.665 \pm 0.020$ \\
Mean $|\text{logit}|$ of kept bits & $14.151 \pm 0.016$ \\
Fraction of helpful flips & $\mathbf{5.3\% \pm 0.5\%}$ \\
\bottomrule
\end{tabular}
\end{table}

The projection flips only 52.5 bits on average (0.16\% of the grid), but
94.7\% of those flips move \emph{away} from the ground truth. The right
inverse finds the minimum-Hamming-weight correction to satisfy the
syndrome, but this algebraically minimal correction has no relationship
to the spatially correct one. The projection operates at the scale of
full rows (32 voxels), while errors are localized at object surfaces
(2--8 voxels). This scale mismatch is the root cause.

% ====================================================================
\section{The Constraint Locality Principle}
\label{sec:locality}

We compare five post-processing strategies applied to the same learned
proposal output (Table~\ref{tab:locality}). Hard projection uses a right
inverse; BP refinement uses the proposal's soft output as variable
beliefs.

\begin{table}[htbp]
\centering
\caption{Clean constraint design ablation. Same learned proposal
(IoU $0.9370\pm0.0028$), different post-processing. 200 ShapeNet validation
samples, spectral level~4, mean $\pm$ SEM over seeds 42/43/44. Extra bits are
additional local syndrome values beyond the fixed global-syndrome payload; this is a
constraint-design ablation, not a rate-matched codec comparison.}
\label{tab:locality}
\begin{tabular}{lccrc}
\toprule
Post-processing & Weight & Overlap & Extra bits & $\Delta$IoU \\
\midrule
None (proposal only) & --- & --- & 0 & baseline \\
Hard projection (global) & 32--1024 & low & 0 & $-0.0588 \pm 0.0058$ \\
\midrule
\multicolumn{5}{l}{\emph{BP refinement with surface constraints:}} \\
\textbf{Surface patches ($2^3$)} & \textbf{8} & \textbf{yes} & \textbf{4{,}205} & $\bm{+0.0363 \pm 0.0026}$ \\
Surface gradient runs & 4 & yes & 12{,}682 & $+0.0603 \pm 0.0031$ \\
Surface neighbor pairs & 2 & yes & 12{,}806 & $+0.0623 \pm 0.0028$ \\
\bottomrule
\end{tabular}
\\[4pt]
\footnotesize{Numbers from \texttt{reproduce\_results.py} Sections~5--6,
clean rsfix artifacts in \texttt{data/clean-v2-may02-rsfix/}.}
\end{table}

The learned decoder gives a smooth nominal rate--distortion diagnostic as more
low-frequency spectral side-information is made available
(Table~\ref{tab:rate-distortion}). These numbers are post-hard-projection
reconstructions before surface-BP refinement; the L4 row matches the
hard-projection setting in Table~\ref{tab:projection}. The table measures the
bandwidth available to the learned proposal plus the cost of the fixed
syndrome, under an 8-bit-per-retained-coefficient accounting model.

\begin{table}[htbp]
\centering
\caption{Clean nominal rate--distortion results on 200 ShapeNet validation
samples, mean $\pm$ SEM over seeds 42/43/44. Bits count the 2{,}977-bit
syndrome, 8 coarse bits, and 8 bits per retained low-frequency spectral
coefficient; no entropy-coded coefficient bitstream is implemented.}
\label{tab:rate-distortion}
\begin{tabular}{lrr}
\toprule
Spectral level & Bits & Mean IoU \\
\midrule
L2 & 3{,}201 & $0.6947 \pm 0.0072$ \\
L3 & 3{,}985 & $0.7941 \pm 0.0081$ \\
L4 & 5{,}729 & $0.8781 \pm 0.0085$ \\
L5 & 8{,}817 & $0.9380 \pm 0.0072$ \\
L6 & 13{,}633 & $0.9729 \pm 0.0036$ \\
L7 & 20{,}561 & $0.9908 \pm 0.0017$ \\
L8 & 29{,}985 & $0.9970 \pm 0.0006$ \\
\bottomrule
\end{tabular}
\end{table}

Because the ShapeNet validation subset is occupancy-sparse, we also stratify
the L4 results by occupied-voxel density (Table~\ref{tab:strata}). The strata
are rank quartiles of the same 200 validation samples, so each row has equal
support. IoU ignores true negatives, so these are not background-accuracy
effects: the table measures overlap on occupied structure.

\begin{table}[htbp]
\centering
\caption{Clean L4 results stratified by occupancy rank quartile. Each stratum
contains 50 validation objects per seed; entries are mean $\pm$ SEM over seeds
42/43/44. Local BP helps in every quartile and helps most in the densest,
hardest objects.}
\label{tab:strata}
\resizebox{\textwidth}{!}{%
\begin{tabular}{@{}llrrrr@{}}
\toprule
Stratum & Occupancy & Raw IoU & Hard $\Delta$ & Patch $\Delta$ & Neighbor $\Delta$ \\
\midrule
Q1 & 0.62--1.68\% & $0.946\pm0.003$ & $-0.059\pm0.007$ & $+0.035\pm0.002$ & $+0.053\pm0.003$ \\
Q2 & 1.68--1.92\% & $0.947\pm0.003$ & $-0.053\pm0.004$ & $+0.034\pm0.002$ & $+0.052\pm0.003$ \\
Q3 & 1.93--2.27\% & $0.944\pm0.003$ & $-0.058\pm0.006$ & $+0.035\pm0.003$ & $+0.056\pm0.003$ \\
Q4 & 2.28--5.12\% & $0.911\pm0.002$ & $-0.065\pm0.006$ & $+0.041\pm0.003$ & $+0.088\pm0.002$ \\
\bottomrule
\end{tabular}
}
\\[4pt]
\footnotesize{Generated by \texttt{stratify\_compression\_results.py};
aggregate JSON: \texttt{data/clean-v2-may02-strata/strata\_summary.json}.}
\end{table}

The Q4 monotonic increase aligns with the locality argument in
Section~\ref{sec:theory}. Denser objects expose more occupied surface, which
in turn places more local checks within the BP support. The single-error
informativeness $P_1(w,p)$ is a per-check quantity, but reconstruction
quality scales with the number of locally informative checks that cover
uncertain voxels. Q4 has the largest fraction of such voxels by
construction, so a fixed local check family yields the largest
$\Delta$IoU there. The same effect is invisible to hard projection because
its high-weight checks have negligible $P_1$ regardless of which stratum
they cover.

All three successful BP variants share three properties:
\begin{enumerate}[nosep]
  \item \textbf{Locality}: few variables per check ($\leq 8$), so
    corrections are spatially contained.
  \item \textbf{Overlap}: adjacent checks share variables, enabling
    BP to propagate information between patches.
  \item \textbf{Targeting}: checks placed where the proposal is uncertain
    (object boundaries), not uniformly over the grid.
\end{enumerate}

Hard projection fails because it operates at the wrong spatial scale
(full rows of 32 voxels) and uses minimum Hamming weight rather than
spatial coherence as its correction objective.

% ====================================================================
\section{Why Constraint Weight Controls Local Informativeness}
\label{sec:theory}

The empirical results in Table~\ref{tab:locality} show that
syndrome-guided correction can either help or hurt a learned proposal. A simple
check-weight calculation explains why high-weight checks are often
uninformative for local correction.

\begin{theorem}[Correction Utility]\label{thm:utility}
Let $\hat{\mathbf{x}}$ be an estimate of $\mathbf{x}^* \in
\mathrm{GF}(q)^n$ where each symbol is independently correct with
probability $p \in (1/2,\,1)$. For a linear check of weight~$w$
(i.e., $\sum_{i=1}^{w} h_i x_i = s$), the probability that the check
contains exactly one error is
\begin{equation}\label{eq:p1}
  P_1(w,\,p) \;=\; w\,(1-p)\,p^{\,w-1}.
\end{equation}
In this single-error regime, the check residual determines the correction
value once the erroneous symbol is identified by the prior or by overlapping
checks.
This function is maximized at $w^* = -1/\!\ln p$ and satisfies
\begin{equation}
  P_1(w,\,p) \;\leq\; w \exp\!\bigl(-(w-1)\ln(1/p)\bigr),
\end{equation}
which decays exponentially for $w \gg w^*$.
\end{theorem}

\begin{proof}
Each of the $w$ symbols in the check is independently wrong with
probability $1-p$. The event ``exactly one wrong'' has probability
$\binom{w}{1}(1-p)^1 p^{w-1}$, giving~\eqref{eq:p1}. Setting
$\partial P_1/\partial w = 0$ (treating $w$ as continuous) yields
$(1-p)p^{w-1}[1 + w\ln p] = 0$, so $w^* = -1/\!\ln p$. The bound
follows from $p^{w-1} = \exp((w{-}1)\ln p)$.
\end{proof}

\smallskip\noindent\textbf{Example.} For a prior with
$p = 0.88$: $w^* \approx 8$, $P_1(4) = 0.33$, $P_1(64) = 5 \times
10^{-4}$. With $m = 32$ checks, local checks ($w{=}4$) yield
${\sim}10.6$ single-error-informative checks; global checks ($w{=}64$) yield
${\sim}0.016$.

\begin{theorem}[Expected utility of residual-only correction]\label{thm:expected}
Consider a weight-$w$ check over $\mathrm{GF}(2)$ on a codeword
$\mathbf{x}^*$, and an estimate
$\hat{\mathbf{x}} = \mathbf{x}^* + \mathbf{e}$ whose entries within the
check support are independently in error with probability $1-p$.
Define the \emph{uniform-residual correction policy}: whenever the
check residual $r \neq 0$, pick a support position $i$ uniformly at
random and flip $\hat{x}_i$. Conditional on $r \neq 0$ (equivalently,
the support contains an odd number of errors), the expected change in
error count is
\begin{equation}\label{eq:expected}
  \mathbb{E}[\Delta\mathrm{err}\mid r\neq 0]
  \;=\; \frac{1}{Z}\!\sum_{\substack{k=1\\ k\,\mathrm{odd}}}^{w}
  \binom{w}{k}(1-p)^{k}\,p^{\,w-k}\,\frac{w-2k}{w},
\end{equation}
with $Z = \sum_{k\,\mathrm{odd}}\binom{w}{k}(1-p)^k p^{\,w-k}$. For
$p > 1/2$ this expectation is strictly positive whenever $w \geq 3$,
with leading order
$\mathbb{E}[\Delta\mathrm{err}\mid r\neq 0] = (w-2)/w + O((1-p)^{2})$
as $p \to 1$.
\end{theorem}

\begin{proof}
Conditional on exactly $k$ errors in the support, the random
correction position lies on an error symbol with probability $k/w$
(the flip removes it: $\Delta\mathrm{err} = -1$) and on a correct
symbol with probability $(w-k)/w$ (the flip introduces a new error:
$\Delta\mathrm{err} = +1$). The conditional expectation is therefore
$(w-2k)/w$. Averaging against the distribution on $k$ given $r \neq 0$
yields~\eqref{eq:expected}. As $p \to 1$, the conditional mass
concentrates at $k=1$, where $(w-2)/w > 0$ for $w \geq 3$.
\end{proof}

\begin{corollary}[A prior is necessary]\label{cor:prior}
No correction policy that selects a support position using only the
residual value can achieve negative expected $\Delta\mathrm{err}$ for
$w \geq 3$ and $p > 1/2$.
\end{corollary}

\begin{proof}
The residual is invariant under permutations of the support, so any
policy depending only on the residual is symmetric across positions
and equivalent in expectation to the uniform-residual policy.
Theorem~\ref{thm:expected} then applies.
\end{proof}

\begin{corollary}[Why high-weight checks are low-utility]\label{cor:global}
For an independent prior with per-symbol accuracy $p$, the probability
that a check is in the unambiguous single-error regime is
$P_1(w,p) = w(1-p)p^{w-1}$. This probability decays exponentially once
$w \gg -1/\ln p$. High-weight checks therefore spend most of their
violations in a regime where the residual alone cannot identify the
correction; the syndrome must be resolved by an external prior or by
overlapping checks, not by the single check.
\end{corollary}

\smallskip
Theorem~\ref{thm:expected} replaces an earlier deterministic claim
that every multi-error correction is harmful, which is false because
multi-error cancellations can sometimes reduce the error count. The
probabilistic version gives the right operational content: without an
external information source --- the learned proposal, or overlapping
local checks --- the residual alone cannot lower the expected error
count for $w \geq 3$ in the high-prior-accuracy regime ($p \to 1$). The
same conclusion holds qualitatively over $\mathrm{GF}(q)$ with $q > 2$
for the analogous policy that pairs a uniform position choice with the
unique residual-cancelling correction at that position. The practical
point is that large checks are rarely locally informative, while small
overlapping checks create many single-error or near-single-error
neighborhoods where the learned proposal and the syndrome can
cooperate.

\section{Content-Adaptive Constraints}
\label{sec:adaptive}

The constraint matrix $\mathbf{H}$ need not be fixed. We derive it
from the spectral side-information $\mathbf{c}$ (low-frequency transform coefficients)
that both encoder and decoder share:

\begin{enumerate}[nosep]
  \item Reconstruct the low-frequency spectral field $\to$ soft spatial proposal.
  \item Detect object surfaces via morphological operations on the
    thresholded proposal.
  \item Place overlapping $2^3$ patch checks at surface locations.
\end{enumerate}

Since both sides perform identical surface detection from the same spectral
side-information, the constraint locations require \textbf{zero additional
signaling bits}. A codec that replaces fixed global checks with local surface
checks would still transmit the local syndrome values themselves
(${\sim}4{,}200$ bits at $32^3$ for patch checks). Table~\ref{tab:locality}
therefore treats these as extra diagnostic bits rather than folding them into
the rate--distortion curve.

% ====================================================================
\section{Learned Proposal and Evaluation Protocol}

The learned proposal is a flat 3D residual network~\cite{resnet}
(width~192, depth~12, 25M~params) with no downsampling. Input:
canonical base solution + spectral field + coarse occupancy map, conditioned
on the syndrome via MLP embedding. Trained on ShapeNet~\cite{3dr2n2}
$32^3$ voxels with random spectral bandwidth per sample.

We note that a small width-48 depth-8 (${\sim}1.68$M~param) model achieves
IoU comparable to the width-192 model on this task --- the learned
proposal is heavily over-parameterized at $32^3$. The constraint system,
not model capacity, is the binding factor for reconstruction quality at
this resolution. We report this only as a qualitative scaling
observation; the multi-seed results in Tables~1--4 use the width-192
model.

Tables~\ref{tab:projection}--\ref{tab:strata} use the width-192
clean v2b models on 200 ShapeNet validation samples, reproduced via
\texttt{reproduce\_results.py}. Full provenance is stored in the local
clean-v2 rsfix and strata data directories. Appendix~\ref{sec:temporal}
reports an auxiliary temporal diagnostic that uses no checkpoint-sensitive
randomness and is identical across the three reproduce jobs.

% ====================================================================
\section{Discussion}

\textbf{When do constraints help?}
Our experiments show that syndrome constraints improve a learned
proposal only in a narrow regime: the successful variants all use
local checks (few variables), overlapping structure (shared variables
between adjacent checks), and surface targeting (placed where the
proposal is uncertain). Hard algebraic projection with global checks
actively degrades quality. This constrains the design space for
syndrome-based spatial coding.

\textbf{The learned proposal does most of the work.}
The clean learned proposal achieves IoU~$0.9370\pm0.0028$ at spectral level~4
before exact hard projection.
Surface BP adds $+0.0363$ to $+0.0623$ depending on the local constraint
family. Post-hoc constraint enforcement is therefore a refinement mechanism,
not the main source of reconstruction quality. Constraint design should be
evaluated jointly with proposal quality: strong learned proposals are the
prerequisite, and local constraints provide targeted refinements.

\textbf{Limitations.}
We tested voxel compression at $32^3$ resolution only. The current
experiments isolate reconstruction behavior but do not yet constitute a
complete entropy-coded production codec. Rate claims must include all
transmitted payloads: fixed syndrome bits, surface-refinement syndromes,
coarse occupancy bits, spectral coefficients, and any variable-rate
selection metadata. The current spectral side field is a controlled
low-frequency diagnostic with nominal 8-bit-per-coefficient accounting, not a
finished quantized coefficient bitstream. Larger resolutions and head-to-head comparison
against established voxel and point-cloud codecs (e.g.,
Draco~\cite{draco}) remain future work.

% ====================================================================
\section{Conclusion}

We studied underdetermined syndrome decoding for 3D occupancy grids with
learned spatial proposals. The main result is not that fixed parity checks
solve compression by themselves; the learned proposal does most of the
reconstruction work. Syndrome constraints help only when they are local,
overlapping, and targeted at uncertain regions of the proposal. This
constraint-locality principle gives a concrete design rule for future
learned syndrome-based occupancy codecs and a clear failure mode for hard
algebraic projection with high-weight checks.

Code and exact reproduction commands are released with the paper family in the
companion portfolio repository.

% ====================================================================
\appendix
\section{Auxiliary Temporal Diagnostic}
\label{sec:temporal}

When multiple frames of the same scene are available, each frame provides an
independent noisy estimate of the true solution. We include this diagnostic
only to show that local BP outputs can be fused over time; it is not part of
the paper's central static-grid contribution. The per-voxel Bayesian
log-odds accumulator is
\begin{equation}
  L_t = \alpha \cdot L_{t-1} + (1-\alpha) \cdot
    \log\frac{\hat{p}_t}{1 - \hat{p}_t},
  \quad L_t \in [-5, 5],
\end{equation}
where $\hat{p}_t$ is the BP-corrected probability at frame $t$,
$\alpha = 0.7$ is the EMA coefficient, and the clamp at $\pm 5$
prevents lock-in. Changed voxels undergo soft decay ($\times 0.3$)
rather than hard reset.

\begin{table}[htbp]
\centering
\caption{Temporal accumulation on a 20-frame morphing scene with
  8\% surface tampering per frame. No neural decoder; uses spectral prior +
  surface BP only. From \texttt{reproduce\_results.py} Section~9. This
  deterministic benchmark is identical across seeds.}
\label{tab:temporal}
\begin{tabular}{lcc}
\toprule
Decoding strategy & Mean IoU & After 5 frames \\
\midrule
Single-frame surface BP & 0.847 & 0.847 \\
\textbf{+ Bayesian accumulation} & \textbf{0.958} & \textbf{0.975} \\
\bottomrule
\end{tabular}
\end{table}

% ====================================================================
\bibliographystyle{plain}
\begin{thebibliography}{12}

\bibitem{gallager}
R.~G.~Gallager.
\newblock Low-density parity-check codes.
\newblock {\em IRE Trans.\ Inf.\ Theory}, 8(1):21--28, 1962.

\bibitem{tanner}
R.~M.~Tanner.
\newblock A recursive approach to low complexity codes.
\newblock {\em IEEE Trans.\ Inf.\ Theory}, 27(5):533--547, 1981.

\bibitem{kfl2001}
F.~R.~Kschischang, B.~J.~Frey, and H.-A.~Loeliger.
\newblock Factor graphs and the sum-product algorithm.
\newblock {\em IEEE Trans.\ Inf.\ Theory}, 47(2):498--519, 2001.

\bibitem{candes2006}
E.~J.~Cand\`{e}s, J.~Romberg, and T.~Tao.
\newblock Robust uncertainty principles: exact signal reconstruction from
  highly incomplete frequency information.
\newblock {\em IEEE Trans.\ Inf.\ Theory}, 52(2):489--509, 2006.

\bibitem{nachmani2016}
E.~Nachmani, Y.~Be'ery, and D.~Burshtein.
\newblock Learning to decode linear codes using deep learning.
\newblock In {\em Allerton Conference}, pp.~341--346, 2016.

\bibitem{nachmani2018}
E.~Nachmani et~al.
\newblock Deep learning methods for improved decoding of linear codes.
\newblock {\em IEEE J.\ Sel.\ Topics Signal Processing}, 12(1):119--131, 2018.

\bibitem{balle2017}
J.~Ball\'{e}, V.~Laparra, and E.~P.~Simoncelli.
\newblock End-to-end optimized image compression.
\newblock In {\em ICLR}, 2017.

\bibitem{minnen2018}
D.~Minnen, J.~Ball\'{e}, and G.~D.~Toderici.
\newblock Joint autoregressive and hierarchical priors for learned image
  compression.
\newblock In {\em NeurIPS}, pp.~10771--10780, 2018.

\bibitem{resnet}
K.~He, X.~Zhang, S.~Ren, and J.~Sun.
\newblock Deep residual learning for image recognition.
\newblock In {\em CVPR}, pp.~770--778, 2016.

\bibitem{3dr2n2}
C.~B.~Choy, D.~Xu, J.~Gwak, K.~Chen, and S.~Savarese.
\newblock {3D-R2N2}: A unified approach for single and multi-view 3D object
  reconstruction.
\newblock In {\em ECCV}, pp.~628--644, 2016.

\bibitem{draco}
Google.
\newblock {Draco}: 3D data compression library.
\newblock \url{https://github.com/google/draco}, 2017.

\bibitem{wynerziv}
A.~D.~Wyner and J.~Ziv.
\newblock The rate-distortion function for source coding with side information
  at the decoder.
\newblock {\em IEEE Trans.\ Inf.\ Theory}, 22(1):1--10, 1976.

\bibitem{slepianwolf}
D.~Slepian and J.~K.~Wolf.
\newblock Noiseless coding of correlated information sources.
\newblock {\em IEEE Trans.\ Inf.\ Theory}, 19(4):471--480, 1973.

\end{thebibliography}

\end{document}
